% Faithful transcription --- original Latin, transcribed from the ORIGINAL publication in the
% Acta Eruditorum, Leipzig, September 1694, pp. 336--338: Jacob Bernoulli's "Constructio Curvae
% Accessus & Recessus aequabilis, ope rectificationis Curvae cujusdam Algebraicae, addenda nuperae
% solutioni Mensis Junii" --- the September sequel to the June "Solutio Problematis Leibnitiani".
% Scan: Internet Archive, Acta Eruditorum 1694 (id s1id13206610), leaves 0436--0438 (the September
% p. 336, distinct from the July p. 336; the paper starts partway down p. 336, after a medical
% review, and ends on p. 338 where "C. H. Z. Constructio Universalis ..." begins). \origpage uses
% the PRINTED Acta page numbers.
%
% NOTE ON SOURCE: first drafted from the 1744 collected Opera (No. LX, pp. 608--612), now re-collated
% against and conformed to the 1694 Acta. Consequences:
%  * The "No. LX" number and Opera pages 608--612 were the collected edition's; dropped. In the Acta
%    this is a running heading, "Jac. B. Constructio Curvae Accessus & Recessus aequabilis ...".
%  * The lettered footnotes (a)--(f) were Cramer's 1744 apparatus, absent from the Acta (removed).
%  * The Opera's closing paragraph "Caeterum quantitas $z^{m}\,dz : \sqrt{a^{4}\pm z^{4}}$ ..." is
%    NOT in the Acta (the paper ends at "proprietas exsistit.") and is not kept --- an Opera addition.
%  * The figure is Fig. 3 on plate "Tab. VI" of the September Acta (the Opera reset it as Tab. XXIV).
%
% Transcription notes for the reviewer:
%  * Author notation is kept faithfully (HOUSESTYLE R3): doubled letters for squares (aa, zz, xx,
%    yy, uu), the ":" ratio/division sign, "&c." and the "&" ampersand, archaic spelling; typography
%    (long-s, ligatures, small caps incl. "Joh. Bernoulli", the superscript-"o" of "No.") is
%    normalized; the author's asides/value-substitutions use round parentheses. The printed "rectangle"
%    glyph for "rectangulum" and the numeral "4" in "Curva 4 dimensionum" are set as printed (the
%    second mention, p. 338, is spelled "quatuor").
%  * This paper first NAMES the lemniscate (p. 337): "seu complicatae in nodum fasciae, sive lemnisci,
%    d'un nœud de ruban Gallis". The printed mark after "jacentis notae octonarii" is a reclining
%    figure-eight (the lemniscate's own shape), set here as $\infty$ (closest available glyph) ---
%    it denotes the numeral 8 lying down, not the modern "infinity" concept.
\documentclass{article}
\usepackage{readmasters}
\begin{document}

\origpage{336}
\section*{Jac.\ B. Constructio Curvae Accessus \& Recessus aequabilis}
\subsection*{ope rectificationis Curvae cujusdam Algebraicae, addenda nuperae solutioni Mensis Junii}

Triplex praecipue modus habetur construendi Curvas Mechanicas sive Transcendentes. Primus, sed ad
praxin parum idoneus, fit per quadraturas spatiorum curvilineorum. Melior est, qui instituitur per
rectificationes Curvarum Algebraicarum; accuratius enim \& expeditius in praxi rectificari possunt
curvae, ope fili vel catenulae ipsis circumplicatae, quam quadrari spatia. Eodem loco habeo illas
constructiones, quae peraguntur absque ulla rectificatione \& quadratura, per solam descriptionem
Curvae alicujus Mechanicae, cujus puncta, licet non omnia, infinita tamen \& quantumvis proxima
geometrice inveniri possunt, qualis solet esse Logarithmica, \& si quae sunt ejus generis aliae.
Optimus vero modus, sicubi haberi possit, ille est, qui peragitur ope alicujus Curvae, quam natura
ipsa absque arte motu quodam celerrimo \& quasi instantaneo ad nutum Geometrae producit; cum
praecedentes modi requirant curvas, quarum delinea-

\origpage{337}
tio sive per motum continuum sive per plurium punctorum inventionem ab Artifice instituatur,
communiter vel lenta vel impedita nimis exsistit: ita constructiones illas Problematum, quae
Hyperbolae quadraturam vel Logarithmicae descriptionem supponunt, caeteris paribus posthabendas
censeo iis, quae ope Catenariae peraguntur, ceu curvae, quam suspensa catena sponte sua citius
induerit, quam reliquis ipse describendis primam manum admoveris.

\emph{Curva Accessus \& Recessus aequabilis} ob solam proponentis Viri commendationem mereri videtur,
ut de constructionibus secundum omnes tres modos illi adornandis simus solliciti. Constructio primi
modi per quadraturas, quam in nupera solutione insuper habui ob jam dictam rationem, institui potest,
sumendo in \emph{Fig.}~3.\ \emph{Tab.}~VI, rectam $EN$ non, ut ibi, ipsi $Au$ sed $iu$ aequalem,
iterumque curvilineo $AEN$ faciendo aequale rectangulum $AO$; hujus enim latus $AP$ nuperae curvae
$AQ$ adaequabitur, sic ut sumpta ad $AB$ \& nunc inventam $AP$ tertia proportionali $A\alpha$,
habeatur punctum $\alpha$ in optata Curva. Tertii modi Constructio, quae fieret mediante Linea
Elastica $AQ$, quam solam ibi tetigi, sine dubio foret optima, si natura alicubi tensiones viribus
tendentibus simpliciter proportionales effecisset: at quia forte nullus invenitur Elater, qui hanc
tensionis legem praecise observet; nec si reperiatur, illud certo constet; idcirco nec isti fidere
satis tutum; praestatque recurrere ad secundum construendi modum, \& quaerere Curvam Algebraicam,
cujus rectificatione scopum assequamur. Talem ex voto sese sistit Curva 4 dimensionum quae hac
aequatione exprimitur $xx + yy = a\sqrt{xx - yy}$, quaeque circum axem $BG\ (2a)$ constituta formam
refert jacentis notae octonarii $\infty$, seu complicatae in nodum fasciae, sive lemnisci,
\emph{d'un nœud de ruban} Gallis. Hujus enim tum altera medietas Curvae Elasticae $AR\varphi$, tum
super nodo $A$ intervallo indeterminatae $AE$ abscissa medietatis portio minor minori $AQ$, major
majori Elasticae portioni $\varphi RQ$ aequatur. Unde quoque ratio constat, illam ad propositi
Problematis constructionem adhibendi.

\rmfigure{figures/fig-3.png}{Fig.~3.}{Fig. 3 (Tab. VI): the elastic curve $AR\varphi$, the lemniscate
looped through the node $A$, the isochronal curve of uniform approach and recession, and the
auxiliary circle, quadrants and construction lines.}

Caeterum sicut infinitae Lineae Mechanicae a Curvae circularis aut parabolicae rectificatione seu
Logarithmicae descriptione dependent; ita quoque praeter Curvam Accessus \& Recessus infinitarum
aliarum, \& partim etiam ipsius Elasticae constructio (ut infra appare-

\origpage{338}
bit) a rectificatione memoratae Algebraicae quatuor dimensionum derivatur. Et ausim asseverare, in
materia constructionis mechanicarum hanc inter caeteras immediate sequi praecedentes; adeo ut
constructio, quae per nullam priorum succedit, proxime vel per hujus Curvae, aut per Lineae
Hyperbolicae aut Ellipticae, aut duarum simul rectificationem tentanda sit. Cujus rei ratio est,
quod post differentialium formulas: $aa\,dz : \sqrt{aa - zz}$, $zz\,dz : \sqrt{aa - zz}$,
$aa\,dz : \sqrt{zz - aa}$, $zz\,dz : \sqrt{zz - aa}$, quae ope quadraturae Circuli \& Hyperbolae
integrantur, simplicissimae fere videantur hae expressiones: $zz\,dz : \sqrt{z^{4}-a^{4}}$,
$aa\,dz : \sqrt{z^{4}-a^{4}}$, $aa\,dz : \sqrt{a^{4}-z^{4}}$, $zz\,dz : \sqrt{a^{4}-z^{4}}$, \&
similes; quarum prima mediante Lineae Hyperbolicae, secunda \& tertia Curvae nostrae lemniscatae,
quarta ejusdem \& Ellipticae rectificatione integratur. Etenim si indeterminata $z$ applicetur ad
Hyperbolam aequilateram (cujus axis transversus est $2a$) ex ipsius centro, satisfaciet arcus
vertici \& applicatae interceptus pro prima formula. Et si eadem $z$ (aut tertia proportionalis ad
$z$ \& $a$, ubi $z$ major quam $a$) ex nodo subtendatur Curvae lemniscatae, inserviet subtensus
arcus pro duabus mediis. Et si denique ex centro Ellipsis (cujus axis minor est $2a$, major
$2a\sqrt{2}$) ipsa $z$ abscindatur in minore, \& per sectionis terminum recta agatur majori
parallela, designabitur intercepta parallelis portio Ellipticae, truncata respondente portione
Curvae lemniscatae, integrale quartae formulae $zz\,dz : \sqrt{a^{4}-z^{4}}$. Unde cum eadem in
\emph{cit.\ fig.} existente $AE = z$, Elementum denotet applicatae $AP$ vel $EQ$ in Elastica, patet
hanc applicatam aequari differentiae duarum portionum Curvae Ellipticae \& Lemniscatae, ideoque
manifestum, quomodo Elastica per rectificationem utriusque construi possit. Et quoniam Lemniscatae
portio, ut supra monui, Elasticae portioni $AQ$ adaequatur, liquet etiam, Ellipticam ipsam adaequari
aggregato $AQ + AP$; quae non inelegans harum Curvarum proprietas exsistit.

\end{document}
